\documentclass[11pt]{article}
\usepackage[margin=1in]{geometry}
\usepackage{amsmath,amssymb,amsthm,booktabs,graphicx,microtype,natbib,hyperref}
\hypersetup{colorlinks=true,linkcolor=blue,citecolor=blue,urlcolor=blue}
\newtheorem{definition}{Definition}
\newtheorem{proposition}{Proposition}
\newtheorem{theorem}{Theorem}
\newcommand{\FinPlan}{FinPlan\mbox{-}RAG}
\title{Point-in-Time Obligation-Aware Retrieval Planning for Financial Question Answering}
\author{Anonymous authors}
\date{}

\begin{document}
\maketitle

\begin{abstract}
Financial questions are often underspecified as retrieval problems: a single
question can require several entities, fiscal periods and filing types, while
documents become available only at particular times. We introduce \FinPlan, a
non-oracle retrieval planner that represents these requirements as explicit
filing obligations and maintains an auditable evidence state. The planner
targets one unresolved obligation at a time, applies a point-in-time cutoff,
and stops only when the obligation set is closed. This design is motivated by
an identification gap: a relevance score is not, in general, sufficient to
decide whether the evidence needed for a financial answer is complete. We
formalize the gap with an observational-equivalence construction and show that
query budget, temporal availability and counterevidence are separate variables
in the decision problem. In a deterministic controlled benchmark with 20 seeds,
seven missingness/conflict scenarios and 5,600 cases, \FinPlan increases closure
from 0.101 to 0.366 relative to single-shot retrieval and uses 8.87 rather than
10 queries on average. Removing counterevidence increases closure to 0.511 but
reduces counterevidence recall from 0.605 to 0.315 and increases overclaiming.
Removing the cutoff produces future leakage in 50.2\% of cases. These results
are mechanism evidence rather than a real-world QA or state-of-the-art claim.
We release a reproducible implementation and a tool/stdio adapter for common
code agents, including Codex and Claude Code.
\end{abstract}

\section{Introduction}
Financial QA systems must align a question with an entity, a reporting period,
and a filing that was actually available when the question was asked. A system
can retrieve a highly relevant sentence and still fail because it omitted a
second required filing, used a later revision, or ignored a limiting disclosure.
Most RAG pipelines expose a ranked list but leave these obligations implicit.

Our research question is: \emph{does an explicit obligation state improve the
next retrieval and stopping decision under a fixed cutoff and query budget?}
We make three bounded contributions. First, we formalize a financial retrieval
obligation as the tuple $(e,y,f,p)$ for entity, fiscal year, filing type and
fiscal period. Second, we give a compact identification argument showing why a
score-only planner cannot guarantee closure when two worlds share the same
observed score but differ in an unobserved required filing. Third, we implement
and challenge the resulting planner with controlled baselines, component
ablations and a temporal leakage stress test. The maintained repository and
the result hash are part of the evidence boundary; no number is filled in from
an unverified external run.

\section{Related work and research gap}
Retrieval-augmented generation combines a retriever with a reader
\citep{lewis2020rag}. Dense and sparse retrievers optimize relevance or lexical
coverage, while temporal QA benchmarks expose long-range and multi-hop
dependencies \citep{maharana2024locomo,wu2024longmemeval}. These systems provide
useful coverage controls but do not by themselves represent which filings are
still required. Financial QA adds a stronger constraint: the answer must be
supported by information available at a declared cutoff.

The gap is therefore not that existing rankers are inaccurate in every case.
It is that a scalar relevance score is an incomplete state for a stopping rule.
Following the lifecycle-aware separation used by SQCAD, we keep candidate
proposal, evidence qualification and reader generation distinct; FinPlan-RAG
specializes that separation to financial filings and point-in-time obligations.
Our hypothesis is deliberately falsifiable:
\begin{quote}
H1: obligation-aware state planning improves closure and query efficiency under
fixed information and budget, while explicit counterevidence and point-in-time
filtering trade coverage for lower overclaiming and lower future leakage.
\end{quote}

\section{Identification gap and theory}
Let $O(q)=\{(e,y,f,p)\}$ denote the obligations parsed from question $q$ and
let $D_{\le c}$ be documents with availability timestamp no later than cutoff
$c$. A planner observes a history $h_t$ and chooses a query target $a_t$ from
the unresolved obligations. Let $S(h_t)$ be any score-only summary used by a
baseline.

\begin{definition}[Closure]
The evidence state is closed if every obligation with a known document id is
covered by at least one retrieved, point-in-time document and all required
slots have a resolved vote. Closure is a process property; it is not answer
correctness.
\end{definition}

\begin{proposition}[Score insufficiency]
Suppose there exist two latent worlds $w_+$ and $w_-$ with the same observed
history and score $S(h_t)$, but with opposite optimal actions for an unresolved
obligation. For any score-only randomized rule that chooses the positive action
with probability $p$, its worst-world regret is at least
\[
\max\{(1-p)m_+,pm_-\}\ge \frac{m_+m_-}{m_++m_-},
\]
where $m_+$ and $m_-$ are the two action margins.
\end{proposition}

\begin{proof}
In $w_+$, taking the negative action incurs at least $m_+$ and occurs with
probability $1-p$; in $w_-$, taking the positive action incurs at least $m_-$
and occurs with probability $p$. Balancing the two lower bounds gives the
stated harmonic-mean floor.
\end{proof}

\begin{theorem}[Obligation-state value decomposition]
Let $x_t=(U_t,E_t,c)$ contain the unresolved obligation set, evidence state and
cutoff, and let $a_t$ be either a targeted retrieval action or stop. If the
one-step utility is $g(x_t,a_t)$ and the transition/observation kernel is
$K_a(dx_{t+1},do_{t+1}\mid x_t)$, then the finite-horizon action value satisfies
\[
Q_t(x,a)=\mathbb{E}[g(x,a)]+\gamma\,\mathbb{E}_{K_a}V_{t+1}(x'),
\qquad V_t(x)=\max_a Q_t(x,a).
\]
For two actions $a,b$, the lifecycle contrast is
\[
\Delta_t(x)=Q_t(x,a)-Q_t(x,b)=\Delta_t^{\mathrm{immediate}}(x)+
\gamma\Delta_t^{\mathrm{continuation}}(x).
\]
Consequently, a score summary $S(x)$ can support an optimal score-only policy
only if $\Delta_t(x)$ is measurable with respect to $S(x)$; otherwise an
observation-equivalent pair with different signs of $\Delta_t$ exists.
\end{theorem}

\begin{proof}
The first identity is the Bellman recursion obtained by conditioning on the
next state and observation. Subtracting the recursion for $b$ from that for
$a$ gives the contrast decomposition. If $\Delta_t$ is not measurable with
respect to $S$, there are $x_+,x_-$ with equal score and opposite contrast
signs; applying the score-insufficiency proposition yields a strictly positive
worst-world regret for every score-only rule.
\end{proof}

The construction is financial rather than purely semantic: $w_+$ contains a
limiting disclosure or a future revision that changes the required filing,
whereas $w_-$ does not. Both can share the same top-ranked passage. Thus the
failure is at the estimand/state layer, not only a reader error.

\begin{theorem}[Point-in-time safety]
If the planner only admits chunks with $available\_at\le c$ into the evidence
state, then any answer context emitted by the planner contains no future
document. If the cutoff predicate is removed, a future document can be
consumed whenever it has higher retrieval relevance than all visible hits.
\end{theorem}

\begin{proof}
The first claim follows directly from the retrieval predicate. For the second,
construct a case with a later chunk whose relevance is strictly larger than all
chunks available by $c$; the maximum-score selection consumes the later chunk.
\end{proof}

\section{FinPlan-RAG method}
\subsection{Obligation parsing}
The parser maps ``Q2 FY 2024'', ``first half of FY 2024'' and annual language to
normalized period requests. Annual Q4 requests map to a 10-K/FY obligation;
quarterly requests map to 10-Q/Q1--Q3. The resolver intersects these signatures
with visible filing metadata and emits stable, sorted obligations.

\subsection{Obligation-aware retrieval}
Given question $q$, obligation $o=(e,y,f,p)$, and unresolved set $U_t$, the
planner forms a targeted query $q\oplus o$ and retrieves one BM25 chunk with
required terms $e$ and $y$. A chunk is eligible only if its timestamp is no
later than the cutoff. The state returned to the reader is
\[
E_t=(O(q),\;C_t,\;n_t,\;closed_t,\;F_t),
\]
where $C_t$ is the covered document set, $n_t$ is query count, $closed_t$ is
the closure flag, and $F_t$ records future ids observed by diagnostics. The
reader is allowed to answer only when $closed_t=1$; otherwise it can abstain or
request more retrieval.

\subsection{Complexity and non-oracle contract}
For $N$ chunks, BM25 indexing is $O(N)$ in the stored term statistics and each
query scans the in-memory chunks. Planning is $O(|O(q)|)$ retrieval steps up to
the budget. The planner receives no hidden document ids, answer labels or
future gold state. This contract is tested by source inspection and deterministic
unit tests.

\section{Experiments}
\subsection{Protocol}
We ran the repository benchmark with 20 seeds, 40 cases per scenario, seven
scenarios (balanced, multi-path, counterevidence-dense, sparse, source-conflict,
future-revision and random-world), and a shared budget of 10. Every method saw
the same paths, documents, cutoff and resolver. The utility was fixed before
evaluation: correct answer $=1$, abstention $=-0.15$, other error $=-1$,
overclaim penalty $=-1$, and query cost $=-0.03$. The full run is identified by
SHA-256 \texttt{E5D207A2...42CAA28}; its compact summary is committed in
\texttt{paper/results/controlled\_summary.json}.

\subsection{Main result and baselines}
Table~\ref{tab:main} reports means across all case-level records. \texttt{single\_shot}
uses the highest-relevance document, \texttt{fixed\_decomposition} cycles through a
static slot list, and \texttt{generic\_adaptive} selects the least-confident slot
without risk weights. \texttt{finplan\_state} is the proposed policy.

\begin{table}[t]
\centering\small
\caption{Controlled mechanism benchmark. Higher is better except overclaim,
future leakage and queries. Values are means over 5,600 cases.}
\label{tab:main}
\begin{tabular}{lrrrrrrr}
\toprule
Method & Acc. & Closure & Overclaim & Counter recall & Queries & Utility & Future leak\\
\midrule
Single-shot & .066 & .101 & .004 & .651 & 10.00 & -.407 & .000\\
Fixed decomposition & .175 & .293 & .014 & .658 & 10.00 & -.363 & .000\\
Generic adaptive & .193 & .343 & .015 & .632 & 10.00 & -.371 & .000\\
FinPlan state & .191 & .366 & .013 & .605 & 8.87 & -.358 & .000\\
\bottomrule
\end{tabular}
\end{table}

Relative to single-shot, FinPlan increases closure by 0.2645 and reduces query
count by 1.1279. Accuracy is similar to generic adaptive (difference -0.0020),
which is an important negative result: obligation state is primarily a process
and stopping improvement in this benchmark, not a universal answer-accuracy
advantage.

\subsection{Ablations and stress tests}
Removing counterevidence raises closure (.511) and lowers queries (7.92), but
counterevidence recall falls from .605 to .315 and overclaim rises from .013 to
.035. Removing adaptive stopping uses all 10 queries and gives utility -.370.
Removing temporal filtering creates future leakage in .502 of cases and gives
utility -.438. These paired changes isolate the roles of counterevidence,
stopping and point-in-time filtering without changing the candidate stream.

\begin{table}[t]
\centering\small
\caption{Component ablations under the same 5,600-case contract.}
\label{tab:ablation}
\begin{tabular}{lrrrrrr}
\toprule
Variant & Closure & Overclaim & Counter recall & Queries & Utility & Future leak\\
\midrule
FinPlan state & .366 & .013 & .605 & 8.87 & -.358 & .000\\
No counterevidence & .511 & .035 & .315 & 7.92 & -.284 & .000\\
No adaptive stop & .337 & .012 & .606 & 10.00 & -.370 & .000\\
No point-in-time & .433 & .021 & .578 & 8.57 & -.438 & .502\\
\bottomrule
\end{tabular}
\end{table}

\subsection{Budget sensitivity}
The quality--cost frontier was checked at budgets $1,2,4,6,10,16$ with the
same 20 seeds and seven scenarios. FinPlan closure increased monotonically from
0.000, 0.000, 0.029, 0.182, 0.366 to 0.455, while counterevidence recall rose
from 0.315 to 0.633. Declared utility decreased from -0.180 to -0.451 because
the fixed query charge eventually dominated the gain from additional closure.
Thus budget 10 is a reproducible operating point, not a universal optimum; the
full sensitivity artifact and per-run hashes are in
\texttt{paper/results/budget\_sensitivity.json}.

\section{Discussion}
The results support H1 only in a bounded sense. Explicit state makes unresolved
obligations visible and permits early stopping, while counterevidence and time
filters impose measurable costs. The no-counterevidence ablation is a useful
challenge to a simplistic ``more closure is better'' interpretation. The
controlled benchmark does not estimate financial answer quality on SEC filings,
does not include a learned reader, and does not establish superiority over
stronger adaptive or dense retrievers. Closure should therefore be reported
alongside correctness, leakage and cost in any future real-data study.

\section{Reproducibility and agent interfaces}
The package requires Python 3.10 and NumPy; `python -m pytest -q` runs the
contract tests. `finplan-rag-agent --schema` emits an OpenAI-compatible tool
schema. The default command reads one JSON request per line and emits one JSON
response per line, so a Codex or Claude Code tool wrapper can call it without
vendor-specific state. The response includes obligations, covered document ids,
selected evidence chunks, closure, future ids, query count and a conservative recommendation to answer
or abstain. It never fabricates a reader answer.

\section{Conclusion}
FinPlan-RAG treats financial retrieval as obligation completion under a cutoff,
not only ranking. The theory identifies when a score-only stopping rule is
insufficient; the controlled evidence shows the measurable trade-off introduced
by explicit counterevidence and temporal filtering. The next decisive step is a
frozen, independently scored financial filing benchmark with a shared reader,
strong retrieval baselines and held-out point-in-time splits.

\appendix
\section{Implementation contract}
The public data structures are \texttt{DocumentObligation}, \texttt{Chunk}, \texttt{RetrievalResult}
and \texttt{EvidenceState}. \texttt{EvidenceState.closed} is true only when every obligation
with a known document id is covered. \texttt{future\_document\_ids} is diagnostic and is
kept separate from the consumed point-in-time context. Reader output is checked
by deterministic numeric support and refusal parsing.

\section{Evidence boundary}
The summary values in Table~\ref{tab:main} trace to the local JSON source hash
listed above. They are not re-derived from prose, and no external SQCAD result
is presented as a FinPlan-RAG result. Real-data results, model weights and large
corpora remain external under the repository storage contract.

\begin{thebibliography}{9}
\bibitem[Lewis et~al.(2020)]{lewis2020rag} Lewis, P. et~al. Retrieval-augmented generation for knowledge-intensive NLP tasks. \emph{NeurIPS} (2020).
\bibitem[Maharana et~al.(2024)]{maharana2024locomo} Maharana, A. et~al. LoCoMo: Long-context conversation memory benchmark. \emph{ACL} (2024).
\bibitem[Wu et~al.(2024)]{wu2024longmemeval} Wu, L. et~al. LongMemEval: Benchmarking long-term memory of large language models. arXiv (2024).
\end{thebibliography}
\end{document}
